% GENERATED FILE. Edit canonical lessons, then run npm run generate:books.
% canonical-source-hash: fnv1a64:9139c837b46f6892
% canonical-lessons: PT-C332-fritar, PT-C332-cozer, PT-C332-assar, PT-C332-grelhar, PT-C332-ferver, PT-R332-revisao-325-332

\chapter{Ao lume}
\label{ch:pt-ao-lume}
\hlchaptermodality{332}

\begin{chapteropening}
\textbf{By the end of this chapter:} \emph{I can say how food was cooked: fried, cooked in water, roasted, grilled or boiled.}

Complete the last lesson of chapter 332: I can say how food was cooked: fried, cooked in water, roasted, grilled or boiled.
\end{chapteropening}

\section[fritar]{fritar — to fry}

\label{lesson:PT-C332-fritar}



\begin{quote}
Before the new one: say the Portuguese for the onion and olive oil base, then the Portuguese for low heat.
\end{quote}

\subsection*{You'll want to know: fritar}
\textbf{fritar} — \textquotedblleft{}to fry\textquotedblright{}.

\begin{quote}
Ela fritou dois ovos para o pequeno-almoço.
\end{quote}

\begin{quote}
\emph{She fried two eggs for breakfast.}
\end{quote}

\emph{Batatas fritas} — chips, crisps.

\begin{cousinweb}
\emph{fritar} is built on \emph{frito} \textquotedblleft{}fried\textquotedblright{} (EtymDB), which goes back to Latin \emph{frīctus}, from \emph{frīgere} \textquotedblleft{}to roast, to fry\textquotedblright{}.
\end{cousinweb}

\begin{grammarlens}[title={What you've built}]
One more word for telling, explaining and describing at B1.
\end{grammarlens}

\subsection*{Your turn}
\begin{itemize}
\raggedright
  \item \emph{Say it:} \emph{fritar}
  \item \emph{Say it:} \emph{fritar}, once more
  \item \emph{Say it:} a sentence of your own with \emph{fritar}
  \item \emph{Recall:} say the Portuguese for the onion and olive oil base, then the Portuguese for low heat, then say \emph{fritar} again
\end{itemize}

\begin{tcolorbox}[breakable,colback=teal!4,colframe=teal!35!black,title={Before you move on}]
What does fritar mean? (\textquotedblleft{}To fry\textquotedblright{}.) Say it once more.
\end{tcolorbox}
\section[cozer]{cozer — to cook in water, to boil (food)}

\label{lesson:PT-C332-cozer}



\begin{quote}
Before the new one: say the Portuguese for low heat, then the Portuguese for to fry.
\end{quote}

\subsection*{You'll want to know: cozer}
\textbf{cozer} — \textquotedblleft{}to cook in water, to boil (food)\textquotedblright{}.

\begin{quote}
Cozemos o arroz em água com sal.
\end{quote}

\begin{quote}
\emph{We cook the rice in salted water.}
\end{quote}

Not \emph{coser}, to sew; they sound the same.

\begin{cousinweb}
\emph{cozer} is inherited from Vulgar Latin \emph{cocere}, from Latin \emph{coquere} \textquotedblleft{}to cook\textquotedblright{} (EtymDB).
\end{cousinweb}

\begin{grammarlens}[title={What you've built}]
One more word for telling, explaining and describing at B1.
\end{grammarlens}

\subsection*{Your turn}
\begin{itemize}
\raggedright
  \item \emph{Say it:} \emph{cozer}
  \item \emph{Say it:} \emph{cozer}, once more
  \item \emph{Say it:} a sentence of your own with \emph{cozer}
  \item \emph{Recall:} say the Portuguese for low heat, then the Portuguese for to fry, then say \emph{cozer} again
\end{itemize}

\begin{tcolorbox}[breakable,colback=teal!4,colframe=teal!35!black,title={Before you move on}]
What does cozer mean? (\textquotedblleft{}To cook in water, to boil (food)\textquotedblright{}.) Say it once more.
\end{tcolorbox}
\section[assar]{assar — to roast, to bake}

\label{lesson:PT-C332-assar}



\begin{quote}
Before the new one: say the Portuguese for to fry, then the Portuguese for to cook in water.
\end{quote}

\subsection*{You'll want to know: assar}
\textbf{assar} — \textquotedblleft{}to roast, to bake\textquotedblright{}.

\begin{quote}
Ao domingo, o meu pai assa um frango.
\end{quote}

\begin{quote}
\emph{On Sundays my father roasts a chicken.}
\end{quote}

\emph{Frango assado} — roast chicken.

\begin{cousinweb}
\emph{assar} is inherited from Latin \emph{assāre} \textquotedblleft{}to roast\textquotedblright{} (EtymDB), from \emph{assus} \textquotedblleft{}roasted, dry\textquotedblright{}.
\end{cousinweb}

\begin{grammarlens}[title={What you've built}]
One more word for telling, explaining and describing at B1.
\end{grammarlens}

\subsection*{Your turn}
\begin{itemize}
\raggedright
  \item \emph{Say it:} \emph{assar}
  \item \emph{Say it:} \emph{assar}, once more
  \item \emph{Say it:} a sentence of your own with \emph{assar}
  \item \emph{Recall:} say the Portuguese for to fry, then the Portuguese for to cook in water, then say \emph{assar} again
\end{itemize}

\begin{tcolorbox}[breakable,colback=teal!4,colframe=teal!35!black,title={Before you move on}]
What does assar mean? (\textquotedblleft{}To roast, to bake\textquotedblright{}.) Say it once more.
\end{tcolorbox}
\section[grelhar]{grelhar — to grill}

\label{lesson:PT-C332-grelhar}



\begin{quote}
Before the new one: say the Portuguese for to cook in water, then the Portuguese for to roast.
\end{quote}

\subsection*{You'll want to know: grelhar}
\textbf{grelhar} — \textquotedblleft{}to grill\textquotedblright{}.

\begin{quote}
No verão, grelhamos peixe no quintal.
\end{quote}

\begin{quote}
\emph{In summer we grill fish in the back yard.}
\end{quote}

\emph{Peixe grelhado} — grilled fish.

\begin{cousinweb}
\emph{grelhar} is built on \emph{grelha} \textquotedblleft{}a grill\textquotedblright{}, borrowed from Old French \emph{greille} (EtymDB), which goes back to Latin \emph{crāticula} \textquotedblleft{}a small gridiron\textquotedblright{}, from \emph{crātis} \textquotedblleft{}a wickerwork\textquotedblright{} (Lewis \& Short).
\end{cousinweb}

\begin{grammarlens}[title={What you've built}]
One more word for telling, explaining and describing at B1.
\end{grammarlens}

\subsection*{Your turn}
\begin{itemize}
\raggedright
  \item \emph{Say it:} \emph{grelhar}
  \item \emph{Say it:} \emph{grelhar}, once more
  \item \emph{Say it:} a sentence of your own with \emph{grelhar}
  \item \emph{Recall:} say the Portuguese for to cook in water, then the Portuguese for to roast, then say \emph{grelhar} again
\end{itemize}

\begin{tcolorbox}[breakable,colback=teal!4,colframe=teal!35!black,title={Before you move on}]
What does grelhar mean? (\textquotedblleft{}To grill\textquotedblright{}.) Say it once more.
\end{tcolorbox}
\section[ferver]{ferver — to boil}

\label{lesson:PT-C332-ferver}



\begin{quote}
Before the new one: say the Portuguese for to roast, then the Portuguese for to grill.
\end{quote}

\subsection*{You'll want to know: ferver}
\textbf{ferver} — \textquotedblleft{}to boil\textquotedblright{}.

\begin{quote}
A água já está a ferver.
\end{quote}

\begin{quote}
\emph{The water is already boiling.}
\end{quote}

The water \emph{ferve}; the food \emph{coze}.

\begin{cousinweb}
\emph{ferver} is inherited from Latin \emph{fervēre} \textquotedblleft{}to boil, to be hot\textquotedblright{} (EtymDB).
\end{cousinweb}

\begin{grammarlens}[title={What you've built}]
One more word for telling, explaining and describing at B1.
\end{grammarlens}

\subsection*{Your turn}
\begin{itemize}
\raggedright
  \item \emph{Say it:} \emph{ferver}
  \item \emph{Say it:} \emph{ferver}, once more
  \item \emph{Say it:} a sentence of your own with \emph{ferver}
  \item \emph{Recall:} say the Portuguese for to roast, then the Portuguese for to grill, then say \emph{ferver} again
\end{itemize}

\begin{tcolorbox}[breakable,colback=teal!4,colframe=teal!35!black,title={Before you move on}]
What does ferver mean? (\textquotedblleft{}To boil\textquotedblright{}.) Say it once more.
\end{tcolorbox}
\section[(dialogue)]{Revisão dos capítulos 325–332}

\label{lesson:PT-R332-revisao-325-332}



\begin{quote}
Say the Portuguese for the character and the protagonist.
\end{quote}

\subsection*{Your turn}
Say these aloud:

\begin{itemize}
\raggedright
  \item the character, the protagonist, the hero, the villain, the ending
  \item the legend, legendary, the myth, the fable, the fairy
  \item the king, the queen, the kingdom, the throne, the crown
  \item the navigator, the caravel, to be shipwrecked, the shipwreck, the treasure
  \item the ancestor, the battle, to invade, to conquer, the revolution
  \item to peel, to chop finely, to grate, to shred, the bowl
  \item to season, the seasoning, to sauté gently, the onion and olive oil base, low heat
  \item to fry, to cook in water, to roast, to grill, to boil
\end{itemize}

\begin{tcolorbox}[breakable,colback=teal!4,colframe=teal!35!black,title={Before you move on}]
The character? (\textbf{a personagem}.) The protagonist? (\textbf{o protagonista}.) The hero? (\textbf{o herói}.) The villain? (\textbf{o vilão}.) The ending? (\textbf{o desfecho}.) The legend? (\textbf{a lenda}.) Legendary? (\textbf{lendário}.) The myth? (\textbf{o mito}.) The fable? (\textbf{a fábula}.) The fairy? (\textbf{a fada}.) The king? (\textbf{o rei}.) The queen? (\textbf{a rainha}.) The kingdom? (\textbf{o reino}.) The throne? (\textbf{o trono}.) The crown? (\textbf{a coroa}.) The navigator? (\textbf{o navegador}.) The caravel? (\textbf{a caravela}.) To be shipwrecked? (\textbf{naufragar}.) The shipwreck? (\textbf{o naufrágio}.) The treasure? (\textbf{o tesouro}.) The ancestor? (\textbf{o antepassado}.) The battle? (\textbf{a batalha}.) To invade? (\textbf{invadir}.) To conquer? (\textbf{conquistar}.) The revolution? (\textbf{a revolução}.) To peel? (\textbf{descascar}.) To chop finely? (\textbf{picar}.) To grate? (\textbf{ralar}.) To shred? (\textbf{desfiar}.) The bowl? (\textbf{a tigela}.) To season? (\textbf{temperar}.) The seasoning? (\textbf{o tempero}.) To sauté gently? (\textbf{refogar}.) The onion and olive oil base? (\textbf{o refogado}.) Low heat? (\textbf{o lume brando}.) To fry? (\textbf{fritar}.) To cook in water? (\textbf{cozer}.) To roast? (\textbf{assar}.) To grill? (\textbf{grelhar}.) To boil? (\textbf{ferver}.)
\end{tcolorbox}
